% ============================================================
%  CLASE 6 — ESTUDIO COMPLETO DE CURVAS
%  Minicurso de Derivadas · Clase de 2 horas
% ============================================================
\documentclass[aspectratio=169]{beamer}
\usepackage{mathwizards-beamer}
\mwbeamer{Clase 6}{Estudio Completo de Curvas}{Minicurso de Derivadas}

\begin{document}

\begin{frame}[plain]
  \titlepage
\end{frame}

% ════════════════════════════════════════════════════════════
\section{Marco Teórico}
% ════════════════════════════════════════════════════════════

\begin{frame}{Primeros Pasos: Dominio, Cortes, Simetría}
  \begin{block}{Dominio}
    Conjunto de $x$ donde $f(x)$ existe. Excluir: denominadores nulos, raíces pares negativas, logaritmos de no positivos.
  \end{block}
  \vspace{6pt}
  \begin{block}{Cortes y simetría}
    \begin{itemize}
      \item Cortes: con $x=0$ $\to$ eje $y$; con $y=0$ $\to$ eje $x$.
      \item $f(-x)=f(x)$ $\to$ par; $f(-x)=-f(x)$ $\to$ impar.
    \end{itemize}
  \end{block}
\end{frame}

\begin{frame}{Asíntotas}
  \begin{block}{Tipos de asíntotas}
    \begin{itemize}
      \item \textbf{Vertical} ($x=a$): $\lim_{x\to a}f(x) = \pm\infty$.
      \item \textbf{Horizontal} ($y=b$): $\lim_{x\to\infty}f(x)=b$ o $\lim_{x\to-\infty}f(x)=b$.
      \item \textbf{Oblicua} ($y=mx+b$): $m=\lim\frac{f(x)}{x}$, $b=\lim(f(x)-mx)$.
    \end{itemize}
  \end{block}
  \vspace{6pt}
  \begin{mwextra}[Ejemplo]
    $f(x)=\frac{x}{x^2-4}$ tiene verticales en $x=\pm2$ y horizontal $y=0$.
  \end{mwextra}
\end{frame}

\begin{frame}{Monotonía, Extremos, Concavidad e Inflexión}
  \begin{block}{Análisis con derivadas}
    \begin{itemize}
      \item Creciente: $f'(x)>0$; decreciente: $f'(x)<0$.
      \item Punto crítico: $f'(x)=0$ o no existe.
      \item Cóncava arriba: $f''(x)>0$; abajo: $f''(x)<0$.
      \item Punto de inflexión: $f''$ cambia de signo.
    \end{itemize}
  \end{block}
  \vspace{6pt}
  \begin{mwpasos}[Pasos]
    \begin{enumerate}
      \item Dominio.
      \item Cortes.
      \item Simetría.
      \item Asíntotas.
      \item Monotonía / extremos ($f'$).
      \item Concavidad / inflexión ($f''$).
      \item Gráfica.
      \item Rango.
    \end{enumerate}
  \end{mwpasos}
\end{frame}

% ════════════════════════════════════════════════════════════
\section{Ejercicios de Clase}
% ════════════════════════════════════════════════════════════

\begin{frame}{Ejercicio 1}
  \begin{mwejercicio}[Ejercicio 1 — Dominio \quad \medio]
    Halle el dominio de $f(x) = \dfrac{x}{x^2 - 4}$.
  \end{mwejercicio}
\end{frame}

\begin{frame}{Ejercicio 2}
  \begin{mwejercicio}[Ejercicio 2 — Cortes y simetría \quad \medio]
    Halle los cortes con los ejes y la simetría de $f(x) = x^3 - x$.
  \end{mwejercicio}
\end{frame}

\begin{frame}{Ejercicio 3}
  \begin{mwejercicio}[Ejercicio 3 — Asíntotas \quad \medio]
    Halle las asíntotas de $f(x) = \dfrac{3x + 1}{2x - 5}$.
  \end{mwejercicio}
\end{frame}

\begin{frame}{Ejercicio 4}
  \begin{mwejercicio}[Ejercicio 4 — Monotonía \quad \dificil]
    Estudie $f(x) = x^3 - 6x^2 + 9x$ (dominio, cortes, simetría, monotonía y extremos).
  \end{mwejercicio}
\end{frame}

\begin{frame}{Ejercicio 5}
  \begin{mwejercicio}[Ejercicio 5 — Concavidad \quad \dificil]
    Estudie la concavidad y los puntos de inflexión de $f(x) = x^4 - 4x^3$.
  \end{mwejercicio}
\end{frame}

\begin{frame}{Ejercicio 6}
  \begin{mwejercicio}[Ejercicio 6 — Extremos relativos \quad \dificil]
    Halle los intervalos de crecimiento y decrecimiento de $f(x) = \dfrac{x - 2}{x^2 - 1}$ y los extremos relativos.
  \end{mwejercicio}
\end{frame}

\begin{frame}{Ejercicio 7}
  \begin{mwejercicio}[Ejercicio 7 — Concavidad e inflexión \quad \dificil]
    Sea $f(x) = 2x^3 + 3x^2 - 12x + 5$. Halle los puntos de inflexión y los intervalos de concavidad.
  \end{mwejercicio}
\end{frame}

\begin{frame}{Ejercicio 8}
  \begin{mwejercicio}[Ejercicio 8 — Estudio completo de racional \quad \muyDificil]
    Estudio completo de $f(x) = \dfrac{x}{x^2 - 4}$.
  \end{mwejercicio}
\end{frame}

\begin{frame}{Ejercicio 9}
  \begin{mwejercicio}[Ejercicio 9 — Estudio con rango \quad \muyDificil]
    Estudio completo de $x^2 y + y - 3x^2 = 0$; determine el rango.
  \end{mwejercicio}
\end{frame}

\begin{frame}{Ejercicio 10}
  \begin{mwejercicio}[Ejercicio 10 — Polinomio completo \quad \muyDificil]
    Estudio completo de $f(x) = (x+2)^2(x-1)^2 + 1$ (gráfica y rango).
  \end{mwejercicio}
\end{frame}

\end{document}
